%==============================================================================
%  ENGR 2405 -- Electrical Circuits I -- Lab 11
%  Simulated 3-Phase AC Circuit
%
%  Compile with pdflatex (run twice so the running header settles).
%  Referenced images : scope_p1_p3.png  scope_p1_p2.png  scope_three_phase.png
%==============================================================================
\documentclass[11pt]{article}

% ---- packages ---------------------------------------------------------------
\usepackage[utf8]{inputenc}
\usepackage[T1]{fontenc}
\usepackage[margin=1in]{geometry}
\usepackage{amsmath,amssymb}
\usepackage{booktabs}
\usepackage{array}
\usepackage{siunitx}
\usepackage{graphicx}
\usepackage{float}
\usepackage{enumitem}
\usepackage{fancyhdr}
\usepackage{caption}
\usepackage{xcolor}
\usepackage{listings}
\usepackage{circuitikz}
\usepackage{tikz}
\usetikzlibrary{arrows.meta,positioning,calc}
\usepackage[hidelinks]{hyperref}

% ---- siunitx v2 / v3 compatibility ------------------------------------------
\ifdefined\qty\else\newcommand{\qty}[2]{\SI{#1}{#2}}\fi
\ifdefined\unit\else\newcommand{\unit}[1]{\si{#1}}\fi

\renewcommand{\arraystretch}{1.15}

% ---- SciLab listing style ---------------------------------------------------
\definecolor{codekw}{rgb}{0.00,0.30,0.70}
\definecolor{codecom}{rgb}{0.20,0.55,0.20}
\definecolor{codestr}{rgb}{0.64,0.08,0.08}

\lstdefinelanguage{Scilab}{
    morekeywords={
        function,endfunction,if,then,else,elseif,end,for,while,do,
        select,case,break,continue,return,clc,clear,disp,mprintf,
        printf,zeros,ones,eye,linspace,inv,det,sum,length,size,abs,
        sqrt,real,imag,exp,sin,cos,tan,atan,plot,plot2d,xlabel,ylabel,
        title,legend,grid,poly,roots,find,pause,error,warning
    },
    sensitive=true,
    morecomment=[l]{//},
    morecomment=[s]{/*}{*/},
    morestring=[b]",
    morestring=[b]',
}
\lstdefinestyle{scilab}{
    language=Scilab,
    basicstyle=\ttfamily\footnotesize,
    keywordstyle=\color{codekw}\bfseries,
    commentstyle=\color{codecom}\itshape,
    numbers=left,
    numberstyle=\tiny\color{gray},
    numbersep=8pt,
    stringstyle=\color{codestr},
    frame=single,
    rulecolor=\color{gray!50},
    breaklines=true,
    showstringspaces=false,
    tabsize=4,
    captionpos=b,
    xleftmargin=14pt,
}
\lstdefinestyle{console}{
    basicstyle=\ttfamily\footnotesize,
    frame=single,
    rulecolor=\color{gray!50},
    breaklines=true,
    tabsize=4,
    captionpos=b,
    xleftmargin=14pt,
}

% ---- running header ---------------------------------------------------------
\pagestyle{fancy}
\fancyhf{}
\lhead{ENGR 2405 -- Lab 11}
\rhead{Abdon Morales}
\cfoot{\thepage}
\renewcommand{\headrulewidth}{0.4pt}

%==============================================================================
\begin{document}

% ----------------------------- title block -----------------------------------
\begin{center}
    {\large\bfseries ENGR 2405 -- Electrical Circuits I}\\[2pt]
    {\large\bfseries Laboratory \#11: Simulated 3-Phase AC Circuit}\\[10pt]
    \begin{tabular}{rl}
        Name:    & Abdon Morales \\
        UT EID:  & am226923 \\
        Date:    & \today \\
    \end{tabular}
\end{center}
\vspace{0.5em}
\hrule
\vspace{1.5em}

% ----------------------------- description -----------------------------------
\section*{Description}
This lab builds a three op-amp network that turns a single sinusoidal source into
a balanced three-phase set. The generator output is phase 1 ($P_1$). A unity
inverter ($U_1$) flips it by \qty{180}{\degree}, an RC section adds another
\qty{60}{\degree} of lag, and a non-inverting amplifier of gain 2 ($U_3$)
restores the amplitude the RC divider took away, producing phase 3 ($P_3$) at
\qty{240}{\degree} of lag. The third op-amp ($U_2$) is an inverting summer that
adds $P_1$ and $P_3$ with equal weight and inverts the result, which lands at
\qty{120}{\degree} of lag and becomes phase 2 ($P_2$). The RC section is the only
frequency-dependent part of the circuit, so a given $R$ and $C$ pair produces a
correct three-phase set at exactly one frequency. We build nine versions of the
circuit using three resistors and three capacitors, calculate the design
frequency for each pair, and then measure the frequency at which the circuit
actually balances along with the phase of $P_2$ and $P_3$ relative to $P_1$.

Throughout the report a phase angle is quoted as a lag with respect to $P_1$, so
$P_2$ at \qty{120}{\degree} lags $P_1$ by \qty{120}{\degree} and $P_3$ at
\qty{240}{\degree} lags it by \qty{240}{\degree}, which is the same signal as a
\qty{120}{\degree} lead. That is the convention the handout uses.

%==============================================================================
\section*{Circuit and Design Equations}
%==============================================================================
The signal chain from $P_1$ to $P_3$ is shown in Figure~\ref{fig:chain}. $U_1$
is an inverting amplifier with $R_1 = R_2 = \qty{10}{\kilo\ohm}$, so its gain is
$-1$ and its output is $P_1$ shifted by \qty{180}{\degree}. That output drives
the series combination of $R_{\text{freq}}$ and $C_{\text{freq}}$, and we tap the
signal across the capacitor:
\[
    \frac{V_C}{V_{U1}} = \frac{1/j\omega C}{R + 1/j\omega C}
    = \frac{1}{1 + j\omega RC},
    \qquad
    \left|\frac{V_C}{V_{U1}}\right| = \frac{1}{\sqrt{1+(\omega RC)^2}},
    \qquad
    \angle = -\tan^{-1}(\omega RC).
\]
Setting the lag to \qty{60}{\degree} gives the design condition
\[
    \tan^{-1}\!\left(\frac{R}{|Z_C|}\right) = \tan^{-1}(\omega RC) = 60^{\circ}
    \quad\Longrightarrow\quad
    \omega RC = \tan 60^{\circ} = \sqrt{3}
    \quad\Longrightarrow\quad
    \boxed{\,f = \frac{\sqrt{3}}{2\pi RC}\,}
\]
At that frequency the divider magnitude is $1/\sqrt{1+3} = 1/2$, which is why
$U_3$ is set for a gain of 2 with $R_9 = R_{10} = \qty{10}{\kilo\ohm}$. The
output of $U_3$ is therefore the same amplitude as $P_1$ and lags it by
$180^{\circ} + 60^{\circ} = 240^{\circ}$.

\begin{figure}[H]
    \centering
    \begin{circuitikz}[american, scale=0.92, transform shape]
        % --- U1 inverting stage ---
        \draw (0,0) node[left]{$P_1$} to[R, l=$R_1$, a=\qty{10}{\kilo\ohm}] (3,0);
        \node [op amp, anchor=-] (U1) at (3,0) {\small $U_1$};
        \draw (U1.+) -- ++(-0.45,0) node[ground]{};
        \draw (3,0) -- ++(0,1.6) coordinate (f1)
              to[R, l=$R_2$, a=\qty{10}{\kilo\ohm}] (f1 -| U1.out) -- (U1.out);
        \draw (U1.out) node[above right,font=\footnotesize]{$180^{\circ}$};
        % --- RC phase shifter ---
        \draw (U1.out) to[R, l=$R_{\text{freq}}$] ++(2.6,0) coordinate (nC);
        \draw (nC) to[C, l_=$C_{\text{freq}}$] ++(0,-2.0) node[ground]{};
        % --- U3 non-inverting x2 ---
        \node [op amp, anchor=+] (U3) at ($(nC)+(1.6,0)$) {\small $U_3$};
        \draw (nC) -- (U3.+);
        \draw (U3.-) -- ++(-0.8,0) coordinate (n3)
              to[R, l_=$R_{10}$, a^=\qty{10}{\kilo\ohm}] ++(0,-1.8) node[ground]{};
        \draw (n3) -- ++(0,1.9) coordinate (f3)
              to[R, l=$R_9$, a=\qty{10}{\kilo\ohm}] (f3 -| U3.out) -- (U3.out);
        \draw (U3.out) -- ++(1.0,0) node[right]{$P_3$}
              node[above left=2pt and -2pt,font=\footnotesize]{$240^{\circ}$};
    \end{circuitikz}
    \caption{Phase-shift chain. $U_1$ inverts, the RC section adds
    \qty{60}{\degree} of lag and halves the amplitude, and $U_3$ restores the
    amplitude with a gain of 2.}
    \label{fig:chain}
\end{figure}

The summing stage is shown in Figure~\ref{fig:summer}. With
$R_5 = R_3 = R_7 = \qty{10}{\kilo\ohm}$ the output is
\[
    P_2 = -\left(\frac{R_7}{R_5}P_1 + \frac{R_7}{R_3}P_3\right) = -(P_1 + P_3).
\]
Writing the two inputs as phasors with $P_1 = V\angle 0^{\circ}$ and
$P_3 = V\angle 120^{\circ}$ (a \qty{240}{\degree} lag is a \qty{120}{\degree}
lead),
\[
    P_1 + P_3 = V\left(1 + \left(-\tfrac{1}{2} + j\tfrac{\sqrt{3}}{2}\right)\right)
    = V\left(\tfrac{1}{2} + j\tfrac{\sqrt{3}}{2}\right) = V\angle 60^{\circ},
\]
so the sum already has the right magnitude, and the inversion built into the
summer moves it to $V\angle 240^{\circ} = V\angle -120^{\circ}$. That is a
\qty{120}{\degree} lag, which is $P_2$. All three outputs then have the same
amplitude and are spaced \qty{120}{\degree} apart, as sketched in
Figure~\ref{fig:phasor}.

\begin{figure}[H]
    \centering
    \begin{circuitikz}[american, scale=0.92, transform shape]
        \draw (0,0) node[left]{$P_1$} to[R, l=$R_5$, a=\qty{10}{\kilo\ohm}] (3,0);
        \draw (0,-1.6) node[left]{$P_3$} to[R, l_=$R_3$, a^=\qty{10}{\kilo\ohm}] (3,-1.6)
              -- (3,0);
        \node [op amp, anchor=-] (U2) at (3,0) {\small $U_2$};
        \draw (U2.+) -- ++(-0.45,0) node[ground]{};
        \draw (3,0) -- ++(0,1.6) coordinate (f2)
              to[R, l=$R_7$, a=\qty{10}{\kilo\ohm}] (f2 -| U2.out) -- (U2.out);
        \draw (U2.out) -- ++(1.0,0) node[right]{$P_2$}
              node[above left=2pt and -2pt,font=\footnotesize]{$120^{\circ}$};
    \end{circuitikz}
    \caption{Summing stage. $U_2$ adds $P_1$ and $P_3$ with unity weight and
    inverts, which places $P_2$ at \qty{120}{\degree} of lag.}
    \label{fig:summer}
\end{figure}

\begin{figure}[H]
    \centering
    \begin{tikzpicture}[scale=1.0, >={Stealth[round]}]
        \draw[gray!40] (0,0) circle (2);
        \draw[gray!60,->] (-2.5,0) -- (2.6,0) node[right,font=\footnotesize]{Re};
        \draw[gray!60,->] (0,-2.5) -- (0,2.6) node[above,font=\footnotesize]{Im};
        \draw[very thick,->] (0,0) -- (0:2)   node[right]{$P_1 = V\angle 0^{\circ}$};
        \draw[very thick,->] (0,0) -- (240:2) node[below left]{$P_2 = V\angle -120^{\circ}$};
        \draw[very thick,->] (0,0) -- (120:2) node[above left]{$P_3 = V\angle +120^{\circ}$};
        \draw[gray] (0.7,0) arc (0:120:0.7);
        \node[font=\footnotesize] at (60:1.0) {$120^{\circ}$};
    \end{tikzpicture}
    \caption{Phasor diagram of the three outputs at the design frequency. Reading
    the sequence as $P_1 \rightarrow P_2 \rightarrow P_3$, each phase lags the
    previous one by \qty{120}{\degree}.}
    \label{fig:phasor}
\end{figure}

%==============================================================================
\section*{Procedure}
%==============================================================================
\begin{enumerate}[leftmargin=*]
    \item We picked three resistors, \qty{1}{\kilo\ohm}, \qty{4.7}{\kilo\ohm}
    and \qty{10}{\kilo\ohm}, and three capacitors, \qty{100}{\nano\farad},
    \qty{220}{\nano\farad} and \qty{470}{\nano\farad}, to use as
    $R_{\text{freq}}$ and $C_{\text{freq}}$. Printed values were used for the
    calculations rather than LCR readings.

    \item We calculated the design frequency $f = \sqrt{3}/(2\pi RC)$ for all
    nine $R$ and $C$ combinations in SciLab before building anything, so we knew
    where to set the generator for each trial.

    \item We built the circuit on the XK700 trainer with
    $R_1 = R_2 = R_3 = R_5 = R_7 = R_9 = R_{10} = \qty{10}{\kilo\ohm}$. The
    supply rails were set to $\pm\qty{12}{\volt}$ from the trainer, and every
    op-amp package got a rail connection before any signal was applied.

    \item We set the function generator for a \qty{5}{\volt} amplitude sine wave
    at the calculated frequency for the first RC pair and applied it at $P_1$.

    \item With channel 1 on $P_1$ and channel 2 on $P_3$, we adjusted the
    generator frequency until the two amplitudes matched, then recorded that
    frequency as $f_{\text{measured}}$.

    \item We measured the phase of $P_3$ with the oscilloscope time cursors. The
    TDS 1002 has no automatic phase readout, so we placed cursors on
    corresponding rising zero crossings of $P_1$ and $P_3$, read $\Delta t$ and
    the period $T$, and computed
    $\phi = 360^{\circ}\,\Delta t / T$.

    \item We moved channel 2 to $P_2$ and repeated the amplitude and phase
    measurement at the same frequency.

    \item We swapped in the next $R$ and $C$ pair and repeated steps 4 through 7
    until all nine combinations were done.
\end{enumerate}

%==============================================================================
\section*{Calculations}
%==============================================================================
The design frequency for each RC pair was computed with the SciLab script below.
The resistor and capacitor arrays hold the values we measured on the LCR meter,
so the printed frequencies are the ones we actually set the generator to.

\begin{lstlisting}[style=scilab, caption={SciLab program \texttt{lab11\_freq.sce}.}]
// ENGR 2405 Lab 11 - design frequency of the RC phase shifter
// The RC tap must lag by 60 deg:  tan(60) = w*R*C = sqrt(3)
//   =>  f = sqrt(3) / (2*pi*R*C)
clear;

R = [1.0e3  4.7e3  10.0e3];     // resistor values, ohms
C = [100e-9 220e-9 470e-9];     // capacitor values, farads

mprintf("   R (ohm)     C (F)      f_calc (Hz)\n");
for i = 1:length(R)
    for j = 1:length(C)
        f = sqrt(3) / (2*%pi*R(i)*C(j));
        mprintf("%10.1f  %10.2e  %11.2f\n", R(i), C(j), f);
    end
end
\end{lstlisting}

\begin{lstlisting}[style=console, caption={SciLab runtime output.}]
   R (ohm)     C (F)      f_calc (Hz)
    1000.0    1.00e-07      2756.64
    1000.0    2.20e-07      1253.02
    1000.0    4.70e-07       586.52
    4700.0    1.00e-07       586.52
    4700.0    2.20e-07       266.60
    4700.0    4.70e-07       124.79
   10000.0    1.00e-07       275.66
   10000.0    2.20e-07       125.30
   10000.0    4.70e-07        58.65
\end{lstlisting}

The prelab table was worked out for \qty{100}{\nano\farad}, \qty{470}{\nano\farad}
and \qty{1}{\micro\farad}. The \qty{1}{\micro\farad} parts were not available at
the bench, so we substituted \qty{220}{\nano\farad} and recomputed the column
above before building anything.

As a check on one row by hand, with $R = \qty{1}{\kilo\ohm}$ and
$C = \qty{100}{\nano\farad}$,
\[
    f = \frac{\sqrt{3}}{2\pi(1000)(100\times10^{-9})}
      = \frac{1.7321}{6.2832\times10^{-4}} = \qty{2756.6}{\hertz},
\]
and back-substituting, $\omega RC = 2\pi(2756.6)(10^{-4}) = 1.7321 = \sqrt{3}$,
so $\tan^{-1}(\sqrt{3}) = 60^{\circ}$ as required.

Each phase angle in the results table was converted from a cursor time
difference using
\[
    \phi = 360^{\circ}\,\frac{\Delta t}{T} = 360^{\circ} f_{\text{meas}}\,\Delta t .
\]
For the same row, $f_{\text{meas}} = \qty{2700}{\hertz}$ gives
$T = \qty{370.4}{\micro\second}$, and the $P_3$ cursor reading
$\Delta t = \qty{244}{\micro\second}$ becomes
$\phi = 360(244)/370.4 = 237.2^{\circ}$.

%==============================================================================
\section*{Data and Results}
%==============================================================================
Table~\ref{tab:results} lists all nine trials. The handout's example table asks
for $\phi(P_1)$, but $P_1$ is the reference for every phase measurement and sits
at $0^{\circ}$ by definition, so we recorded $\phi(P_3)$ in that column
instead.\footnote{$\phi(P_3)$ is what the procedure actually asks to measure, so
we read the column heading as a typo.} Phases are lags with respect to $P_1$;
the design targets are \qty{240}{\degree} for $P_3$ and \qty{120}{\degree} for
$P_2$. The raw cursor readings are included so the angles can be traced back to
what was on the screen.

\begin{table}[H]
    \centering
    \small
    \caption{Measured results for all nine RC combinations. Phase angles are
    computed from the cursor readings as $\phi = 360^{\circ} f_{\text{meas}}\Delta t$.}
    \label{tab:results}
    \begin{tabular}{cc rr r rr rr}
        \toprule
        $R$ & $C$ &
        $f_{\text{calc}}$ & $f_{\text{meas}}$ & \% err &
        $\Delta t_{P3}$ & $\phi(P_3)$ &
        $\Delta t_{P2}$ & $\phi(P_2)$ \\
        (\unit{\kilo\ohm}) & (\unit{\nano\farad}) &
        (\unit{\hertz}) & (\unit{\hertz}) & &
        (\unit{\milli\second}) & (\unit{\degree}) &
        (\unit{\milli\second}) & (\unit{\degree}) \\
        \midrule
        1.0  & 100 & 2756.64 & 2700 & $-2.05$ & 0.244 & 237.2 & 0.124 & 120.5 \\
        1.0  & 220 & 1253.02 & 1200 & $-4.23$ & 0.540 & 233.3 & 0.260 & 112.3 \\
        1.0  & 470 &  586.52 &  586 & $-0.09$ & 1.120 & 236.3 & 0.580 & 122.1 \\
        4.7  & 100 &  586.52 &  586 & $-0.09$ & 1.140 & 240.5 & 0.560 & 118.1 \\
        4.7  & 220 &  266.60 &  266 & $-0.23$ & 2.520 & 241.3 & 1.240 & 118.8 \\
        4.7  & 470 &  124.79 &  125 & $+0.17$ & 5.300 & 238.5 & 2.600 & 117.0 \\
        10.0 & 100 &  275.66 &  276 & $+0.12$ & 2.400 & 238.5 & 1.200 & 119.2 \\
        10.0 & 220 &  125.30 &  125 & $-0.24$ & 5.400 & 243.0 & 2.600 & 117.0 \\
        10.0 & 470 &   58.65 &   59 & $+0.59$ & 11.400 & 242.1 & 5.700 & 121.1 \\
        \bottomrule
    \end{tabular}
\end{table}

Table~\ref{tab:dev} pulls out how far each angle sits from its target. It also
lists the angle an ideal RC section would produce at the frequency we actually
ran at, which separates the part of the phase error that follows from being off
the design frequency from the part that does not.

\begin{table}[H]
    \centering
    \small
    \caption{Phase deviation from the design targets of \qty{240}{\degree} and
    \qty{120}{\degree}. The ideal column is $180^{\circ} + \tan^{-1}(\omega RC)$
    evaluated at $f_{\text{meas}}$, which is what $P_3$ should read given that we
    were not exactly on the design frequency.}
    \label{tab:dev}
    \begin{tabular}{cc r rr rr}
        \toprule
        $R$ & $C$ &
        $\phi(P_3)$ ideal & $\phi(P_3)$ meas. & dev. &
        $\phi(P_2)$ meas. & dev. \\
        (\unit{\kilo\ohm}) & (\unit{\nano\farad}) &
        (\unit{\degree}) & (\unit{\degree}) & (\unit{\degree}) &
        (\unit{\degree}) & (\unit{\degree}) \\
        \midrule
        1.0  & 100 & 239.48 & 237.2 & $-2.8$ & 120.5 & $+0.5$ \\
        1.0  & 220 & 238.92 & 233.3 & $-6.7$ & 112.3 & $-7.7$ \\
        1.0  & 470 & 239.98 & 236.3 & $-3.7$ & 122.1 & $+2.1$ \\
        4.7  & 100 & 239.98 & 240.5 & $+0.5$ & 118.1 & $-1.9$ \\
        4.7  & 220 & 239.94 & 241.3 & $+1.3$ & 118.8 & $-1.2$ \\
        4.7  & 470 & 240.04 & 238.5 & $-1.5$ & 117.0 & $-3.0$ \\
        10.0 & 100 & 240.03 & 238.5 & $-1.5$ & 119.2 & $-0.8$ \\
        10.0 & 220 & 239.94 & 243.0 & $+3.0$ & 117.0 & $-3.0$ \\
        10.0 & 470 & 240.15 & 242.1 & $+2.1$ & 121.1 & $+1.1$ \\
        \midrule
        \multicolumn{4}{r}{Mean deviation\quad} & $-1.0$ & & $-1.5$ \\
        \multicolumn{4}{r}{RMS deviation\quad}  & $3.1$  & & $3.2$ \\
        \bottomrule
    \end{tabular}
\end{table}

%==============================================================================
\section*{Observations \& Discussion}
%==============================================================================
The measured and calculated frequencies agreed closely for seven of the nine
combinations, all within \qty{0.6}{\percent} of the design value and six of them
within \qty{0.25}{\percent}. The two outliers are the two highest-frequency
trials, \qty{1}{\kilo\ohm} with \qty{100}{\nano\farad} ($-2.05\%$) and
\qty{1}{\kilo\ohm} with \qty{220}{\nano\farad} ($-4.23\%$), and in both the
recorded frequency is a round number, \qty{2700}{\hertz} and
\qty{1200}{\hertz}. Those are readout results rather than physical ones. Above
about \qty{1}{\kilo\hertz} we were reading the generator to the nearest
\qty{100}{\hertz}, so the discrepancies of \qty{57}{\hertz} and \qty{53}{\hertz}
are roughly half a step in each case. The remaining seven rows sat where the
display resolution was fine enough to land on the calculated value, and they did.

That column is a weaker test than it looks, and it is worth being honest about
that. We set $f_{\text{meas}}$ by adjusting the generator until the amplitude of
$P_3$ matched $P_1$, and with the gain of 2 from $U_3$ included the balance
condition is
\[
    A(f) = \frac{2}{\sqrt{1 + (\omega RC)^2}}, \qquad
    \omega RC = \sqrt{3}\,\frac{f}{f_0},
\]
which gives $A = 0.964$ at \qty{5}{\percent} above $f_0$ and $A = 0.929$ at
\qty{10}{\percent} above. On a \qty{5}{\volt} amplitude signal that is only
\qty{180}{\milli\volt} and \qty{355}{\milli\volt} of change, so matching the two
traces by eye pins the frequency to something like $\pm\qty{5}{\percent}$. The
phase reading is no sharper on its own: differentiating
$\phi_{RC} = \tan^{-1}(\sqrt{3}f/f_0)$ at $f = f_0$ gives about
\qty{0.25}{\degree} per \qty{1}{\percent} of frequency error against a cursor
resolution of one to two degrees. The sub-percent agreement in the frequency
column is therefore not evidence that the model is accurate to a fraction of a
percent. What the two criteria establish together is a consistency check: the
amplitudes balance and all three phases sit near their targets at the same
frequency, and the model predicts that frequency to within the resolution of
both.

The phases themselves came out well. $\phi(P_3)$ landed within
\qty{3.7}{\degree} of \qty{240}{\degree} in eight of nine trials and $\phi(P_2)$
within \qty{3.0}{\degree} of \qty{120}{\degree} in eight of nine, with RMS
deviations of \qty{3.1}{\degree} and \qty{3.2}{\degree}, about
\qty{1}{\percent} of a full cycle. Both means are slightly negative,
$-1.0^{\circ}$ and $-1.5^{\circ}$, so there is a small systematic shortfall in
lag on top of the scatter. Part of that is not an error at all: if $f < f_0$
then $\omega RC < \sqrt{3}$ and the RC section lags by less than
\qty{60}{\degree}, so $P_3$ should come in under \qty{240}{\degree}.
Table~\ref{tab:dev} evaluates $180^{\circ} + \tan^{-1}(\omega RC)$ at the
frequency we actually used, and for the \qty{1}{\kilo\ohm} and
\qty{220}{\nano\farad} row that ideal value is \qty{238.9}{\degree}, which
accounts for about \qty{1.1}{\degree} of the \qty{6.7}{\degree} gap. That same
row has by far the worst $\phi(P_2)$ at \qty{112.3}{\degree}, which follows,
since $P_2$ is assembled from $P_1$ and $P_3$ by the summer and inherits
whatever error $P_3$ carries. It is the only row where we were meaningfully off
the design frequency and the only row where both angles are more than
\qty{5}{\degree} out.

The rest of the scatter comes from resolution and from the parts. The TDS 1002
has no automatic phase function, so every angle came from placing two cursors on
corresponding zero crossings and reading $\Delta t$ against the period. Our
$\Delta t$ readings carry two or three significant figures, and one unit in the
last figure is already worth about a degree: \qty{1}{\micro\second} out of a
\qty{370}{\micro\second} period is \qty{0.97}{\degree} in the first row, and
\qty{0.1}{\milli\second} out of \qty{16.9}{\milli\second} is \qty{2.1}{\degree}
in the last. That alone brackets most of what we saw. On the component side,
$f_0$ depends on the product $RC$, so a tolerance in either part propagates
one-for-one into the calculated frequency, and \qty{5}{\percent} resistors with
\qty{10}{\percent} capacitors put a worst case band of roughly
\qty{15}{\percent} around each nominal value. We used printed values rather than
LCR readings, so entries in the $f_{\text{calc}}$ column could be off by more
than the differences under discussion. The analysis also treats the capacitor
node as unloaded when it actually drives the non-inverting input of $U_3$, whose
input capacitance and the breadboard together add a few picofarads. That matters
most for the \qty{100}{\nano\farad} trials at the top of our range and is
negligible at \qty{470}{\nano\farad}. Each op-amp stage contributes a small lag
of its own at finite loop gain as well, and those accumulate through the chain,
consistent with the mean deviation in $\phi(P_2)$ being slightly larger in
magnitude than the one in $\phi(P_3)$.

As for where a circuit like this gets used, the point of it is producing a
balanced three-phase reference from a single-phase source and a handful of parts.
On a bench that lets us exercise three-phase equipment such as phase sequence
indicators, protective relays, or the front end of a three-phase rectifier
without needing an actual three-phase supply. In a control setting the same
three signals can serve as the reference for the gate drive of a three-phase
inverter driving an induction motor, where the reference frequency sets the motor
speed, and the same idea appears in signal processing where fixed phase offsets
are generated this way for quadrature modulation. The limitation is the one these
nine trials demonstrate directly: the shift is correct at a single frequency for
a given $R$ and $C$, so a variable-frequency version needs $R$ or $C$ to track
the frequency instead of staying fixed. That is why the practical modern approach
is a digital one, where three phases are read out of a lookup table clocked at
whatever rate the application needs and the \qty{120}{\degree} spacing is exact
by construction.

%==============================================================================
\end{document}
